\documentclass{article}
\usepackage{amsmath, amsfonts, amssymb}
\usepackage{graphicx}
\usepackage{titlesec}
\usepackage{lipsum}
\usepackage[utf8]{inputenc}
\usepackage{xcolor}
\usepackage{listings}
\usepackage{lipsum}
\usepackage{scrextend}
\usepackage{tikz}
\usepackage{hyperref}
\usepackage{color}
\usepackage{fancyvrb}
\usepackage[left=2.5cm,top=3cm,right=2.5cm,bottom=3cm,bindingoffset=0.5cm]{geometry}

\title{Modeling Malcognition}
\author{Luc Alexander L{\"u}thi}
\date{May 2nd, 2026}

\begin{document}
	\maketitle
	\tableofcontents
	\section{Introduction}
	Adversarial systems are often treated as separate domains based on the surface material of the conflict. Malware is treated as a computer security concern, propaganda as a social or psychological concern, debate as a linguistic concern, and cognitive dissonance as a mental concern. The claim of this paper is that these domains share a primitive structural shape. In each case there exists a system with protected invariants, an interface through which signals enter, a set of constraints which supposedly preserve those invariants, and an adversary attempting to produce a signal which causes the system to violate, abandon, or misrecognize its own protected state.\\\\
	To model this shape, we begin from argumentation. Argumentation exposes the basic adversarial structure clearly: a claim is made, a counterexample is offered, support is demanded, and the position is either repaired or abandoned. Security claims have the same form. To say that a system is secure is to say that its constraints are sufficient to preserve some invariant over all relevant inputs. An exploit is therefore not merely a program, but a counterexample to a claimed sufficiency of constraint. This reframes exploitation as an adversarial proof activity over a computational substrate.\\\\
	This paper then generalizes that primitive shape into a substrate agnostic model. A substrate may be computational, cognitive, linguistic, or social, provided that it supports terms, transformations, signals, and effects. Defense constructs constrained transformations which attempt to preserve invariants. Attack searches an enumerable domain for signals or terms which demonstrate constraint failure. Where the system is layered, each demonstrated failure expands visibility or capability, modeling escalation, lateral movement, persuasion, or conceptual collapse.\\\\
	The central concept introduced here is malcognition. Malcognition names the class of adversarial signals which act on belief like structures. A malcognitive signal is not merely false information, nor merely hostile input. It is a signal engineered to be metabolized by the target system in such a way that the system updates, repairs, or preserves itself incorrectly. Its danger lies in the fact that the system may continue to believe that its invariants are intact while the adversary has altered the practical meaning of those invariants.\\\\
	The aim is not to reduce human cognition to ordinary computation, nor to claim that malware and propaganda are identical in their implementation. Rather, the aim is to describe the abstract adversarial skeleton common to both. By representing beliefs, constraints, repairs, and signals as relations in a colored edge hypergraph, we obtain a model in which security failures can be discussed through a shared language of invariant stress.
	\section{Argumentation}
	The most intuitive instantiation of information combat exists in simple logical argumentation, or debate. In debate, claims can be made, challenged, refined, or supported. Any claim can be made, butonly relevant ones are considered, the rest may be dismissed contextually. Claims in linguistic argumentation rely on such massive logical infrastructure that we cannot reliably prove most of them exhaustively, so they are supported up until agreed upon sufficience. Claims may be countered, which prompts the one who made the claim to either refute the counter, make a new claim to replace it, or further support their claim. If claims are positions to be proven, this forms an infinitely grained proof tree.\\\\
	A security context may be anchored in argumentation, to say something is secure is a claim that present constraints are sufficient to preserve a secure invariant. Circumvention of this security is a counter example to that claim.
	\section{The Primitive Shape of Exploitation}
	For any computational substrate which implements a simple interface, you can produce an advesarial game of construction and exploitation. 
The substrate of choice must be computational, in that you must be able to take a term $A$ and naturally derive its evaluation to a term $B$, it must allow terms which take input and produce output, and it must be able to encode, wheher artificially or intrinsically, side effectful computations. Additionally, the substrate must be able to encode an identity term. \\\\
	A player in defense of their position must construct a term $F$ which encodes a series of constraints $C$, and an optional filter term $S$ such that $F \circ S$ encodes all constraints in $C$. Some constraints are designed to not be exhaustively provable, but for the constraints which are, proofs that $F \circ S$ fulfil those constraints are required. A default identity term is chosen for $S$. A player in attack is given an allowed enumerable domain, and is allowed to attempt to prove constraint failures. \\\\
	This constitutes one layer. Multiple layers may be applied to increase the difficulty of the problem, modeling stages of system compromise. Each constraint failure proven via counter example by an attacker must allow visibility into the next layer, and optionally may provide an expanded domain for the player in attack.\\\\
	While this is malleable based on the instantiated substrate, a sensible set of constraint templates follow:
	\begin{itemize}
		\item[] $\exists x$ s.t $(F\circ S)(x) = A$
		\item[] $\neg\exists x$ s.t $(F\circ S)(x) = A$
		\item[] $\forall x, (F\circ S)(x)$ produces side effect A
		\item[] $\forall x, (F\circ S)(x)$ does not produce side effect A
	\end{itemize}
	Constraints act as information gates for subsequent layers and regions of accessible domain for the attacker. When a constraint is broken, the defender has the opportunity to repair by rewriting the term that broke. Multiple terms are allowed per layer with differing constraint sets. When a term is repaired, be it the only broken term in a layer, the domain rerestricts to that layers domain, and the layer visibility is reset to that layer. Some criteria may be given for when a layer is irrepairable, or alternatively when a layer serves as a checkpoint for the attacker beyond which the domain may no longer be rerestricted on repair.\\\\
	It should be noted at this point that constraints are central to any supposed model over a substrate. Any system which can be exploited has a set of constraints which are insufficient to guard against invariant break.
	\section{Malcognition}
	Cognitive dissonance is described as a mental phenomenon in which people unknowingly or subconsciously hold fundamentally conflicting cognitions. This shape exists beyond cognition however. Computing systems, argumentation, and physical security systems all have invariants which they protect. These invariants often have a heirarchy of priority. In cognition this is clearly seen with core identity beliefs, "I am not stupid". This is a belief which under contradiction would likely take priority. Structurally this looks like the cost of updating beliefs being higher for certain beliefs than others. When contradicting beliefs enter the mind and are brought to attention, the beleif structure collapses to preserve the highest priotity invariant, selecting the update that leads to a less costly restructuring. \\\\
	A malcognitive signal is essentially a signal with ulterior motives. A secured system will have central beliefs, it will have belief structure. A malcognitive event either convinces that system of some belief, or culls a secured belief while convincing the system that all its invariants are still held and that the constraints it relies on are enough to secure its invariants. The mechanization of malcognition involves a weaponization of cognitive dissonance, where belief collapse is engineered to produce a certain set of resulting beleifs.\\\\
	The clearest instantiations of malcognitive efforts in cognitive and computational systems respectively are propagnda and malware. Both really take the same shape, as they are malcognitive beliefs engineered to become core beliefs, to take advantage of the system without the system noticing, and to spread. These are distinctly engineering efforts in both cases. Malcognition is methodical, just as in order to make good malware one must know a lot about the operating system and computer it will be running on, producing good propaganda requires very detailed knowledge of human psychology.\\\\
	\subsection{Semantic Substrate Agnostic Modeling}
	Our primitive semantic shape was substrate agnostic but was a husk, it describes the modeling space with no actionable primitives for what a belief is, how a constraint is represented, what it looks like for invariants to be stressed, and no sense of what malcognitive events look like in contrast to normal beliefs. Our refinement will present an abstract representation which remains substrate agnostic but ellicits emergent properties of cognitive and computive systems in regard to adversarial dynamics and exploitability. \\\\
	We begin with a set of relations. These relations compose a colored edge hypergraph. We then proceed with a set of rules which operate on these relations. There are two concerns of these rules: generation, and constraint. Generation rules comprise the API to the system. They read signals and produce new relations from those signals. Constraints ensure that the set of relations does or does not exhibit certain properties, and conditionally mutate the set of relations in an effort to repair broken states and preserve invariants. Our function in this case is the interface of signals, and our filter is the set of constraints.\\\\
\end{document}
